% !TeX root = ../../progress-sem1.tex
% 2.x ทฤษฎีที่เกี่ยวข้อง
\section{ทฤษฎีที่เกี่ยวข้อง}

\subsection{UX/UI Development Life Cycle: UX/UI DLC}
\thaijustify{
  กระบวนการออกแบบประสบการณ์ผู้ใช้และส่วนต่อประสานผู้ใช้เป็นกระบวนการที่เน้นผู้ใช้เป็นศูนย์กลางและต้องอาศัยการทำซ้ำอย่างต่อเนื่อง (Iterative process) โดยทั่วไปนักออกแบบจะต้องผ่านขั้นตอนตั้งแต่การรวบรวมความต้องการ การทำความเข้าใจผู้ใช้ (Empathize) การระบุปัญหา (Define) การระดมความคิด (Ideate) การสร้างต้นแบบ (Prototype) และการทดสอบ (Test)
}
\thaijustify{
  อย่างไรก็ตาม ในทางปฏิบัตินักออกแบบมักเผชิญข้อจำกัดด้านเวลาและทรัพยากร
  ทำให้การสำรวจแนวคิด การรับข้อเสนอแนะ และการปรับปรุงงานซ้ำหลายรอบทำได้ไม่เต็มที่
  Suh et al. \cite{suh2025storyensemble} ชี้ให้เห็นว่าข้อจำกัดดังกล่าวลดโอกาสในการสำรวจทางเลือกและย้อนกลับไปปรับแก้งานในขั้นตอนก่อนหน้า ขณะที่
  Feng et al. \cite{feng2023designing} แสดงให้เห็นว่าการสร้าง wireframe
  ด้วยมือเป็นงานที่ใช้เวลาและแรงงาน ซึ่งสามารถใช้ LLM ช่วยสร้าง wireframe
  จากความต้องการทางการออกแบบได้ และ Yuan et al. \cite{prototypeflow_humanai_ui_design} พบว่าการสร้าง
  UI ด้วย Generative AI ควรเปิดโอกาสให้นักออกแบบตรวจสอบและปรับผลลัพธ์ระหว่างกระบวนการ
  แทนการสร้างผลลัพธ์เพียงครั้งเดียว
}

\thaijustify{จากงานเหล่านี้จึงเห็นได้ว่า Generative AI สามารถช่วยลดภาระของงานที่ต้องทำ
  ด้วยมือและสนับสนุนกระบวนการออกแบบแบบวนซ้ำ ทำให้นักออกแบบมีเวลาสำหรับการทดสอบกับผู้ใช้และนำข้อเสนอแนะกลับมาปรับปรุงต้นแบบได้มากขึ้น
}

\subsection{ระเบียบวิธีการทดสอบการใช้งาน (Usability test)}
\thaijustify{
  การทดสอบกับผู้ใช้เป็นหัวใจสำคัญในการหาข้อบกพร่องและข้อควรปรับปรุงของงานออกแบบ
  โดยงานวิจัยชี้ให้เห็นว่าการทดสอบบ่อยครั้งกับกลุ่มตัวอย่างขนาดเล็กให้ข้อมูลเชิงลึกที่นำไปปรับปรุงได้จริงมากกว่าการทดสอบสเกลใหญ่แต่นาน ๆ ครั้ง โดย Nielsen \cite{nielsen2000five} เสนอว่า
  การทดสอบเชิงคุณภาพสามารถเริ่มต้นด้วยผู้เข้าร่วมประมาณ 5 คนต่อรอบ เนื่องจากภายใต้
  แบบจำลองการค้นพบปัญหาด้านการใช้งาน ผู้เข้าร่วม 5 คนสามารถค้นพบปัญหาได้ประมาณ
  ร้อยละ 85 และการเพิ่มผู้เข้าร่วมหลังจากนั้นมีแนวโน้มให้ข้อมูลใหม่ลดลงด้วย จึงควรนำผลที่ได้
  ไปปรับปรุงต้นแบบและทดสอบซ้ำในรอบถัดไป
  แนวคิดนี้สอดคล้องกับ Virzi \cite{virzi1992refining} ซึ่งพบว่าผู้เข้าร่วมประมาณ 4-5 คนสามารถค้นพบปัญหา
  ด้านการใช้งานได้ประมาณร้อยละ 80  และแบบจำลองของ
  Nielsen และ Landauer \cite{nielsen1993mathematical} ที่แสดงให้เห็นว่าประโยชน์จากการเพิ่มจำนวนผู้เข้าร่วมจะลดลง
  เรื่อย ๆ เมื่อปัญหาส่วนใหญ่ถูกค้นพบแล้ว อย่างไรก็ตาม จำนวนดังกล่าวเหมาะสำหรับการทดสอบเชิงคุณภาพเพื่อค้นหาปัญหาและปรับปรุงงานแบบวนซ้ำ ไม่ใช่เกณฑ์ตายตัวสำหรับการทดสอบทุกประเภท ใน UX/UI DLC มีระเบียบวิธีทดสอบการใช้งานที่นิยมใช้ ดังนี้
  \begin{itemize}
    \item \textbf{Usability Test} เป็นการทดสอบเพื่อให้ผู้เข้าร่วมทดสอบลองใช้งานระบบและหาจุดที่ใช้งานยากหรือติดขัด ซึ่งในขณะที่ทำการทดสอบผู้เข้าร่วมทดสอบจำเป็นต้องคิดไปด้วยพูดไปด้วย (Thinking-Aloud) โดยนักออกแบบจะคอยสังเกตพฤติกรรมและรับฟังความคิดเห็นของผู้เข้าร่วมทดสอบขณะทำภารกิจที่กำหนด \cite{rubin2008handbook, dumas1999practical, nielsen1993usability}
    \item \textbf{A/B Test} เป็นกระบวนการนำเสนอทางเลือกการออกแบบ (Variants) ที่แตกต่างกันให้กับผู้ใช้ เพื่อเปรียบเทียบประสิทธิภาพการใช้งานและรวบรวมข้อเสนอแนะก่อนที่จะตัดสินใจออกแบบในขั้นตอนถัดไป ในภาคอุตสาหกรรม A/B Test มักหมายถึงการทดลองแบบควบคุมบนระบบจริง (Online Controlled Experiment) ที่วัดผลด้วยสถิติโดยเมทริกซ์ต่าง ๆ เช่น conversion rate, click-throughs rate, task completion rate จากผู้ใช้จำนวนมาก \cite{kohavi2009controlled, kohavi2020trustworthy} อย่างไรก็ตาม วิธีนี้ให้ผลเป็นตัวเลขว่าทางเลือกใดทำงานได้ดีกว่า แต่ไม่อธิบายเหตุผลเบื้องหลังพฤติกรรมของผู้ใช้ จึงควรใช้ร่วมกับการรวบรวมข้อมูลเชิงคุณภาพ \cite{nielsen2005abtesting} % ในโครงงานนี้ A/B Test หมายถึงการให้ผู้ใช้กลุ่มเล็กลองใช้ต้นแบบสองรูปแบบและให้ข้อเสนอแนะในแต่ละรอบการทดสอบ ซึ่งเป็นแนวทางเชิงคุณภาพ
    \item \textbf{Comparative Usability Test} มีหลักการคล้าย A/B Test แต่เน้นการทดสอบเชิงคุณภาพและให้ผู้ใช้เปรียบเทียบดีไซน์ตั้งแต่สองรูปแบบขึ้นไปที่มีเป้าประสงค์เดียวกัน งานวิจัยพบว่าการให้ผู้ใช้เห็นทางเลือกการออกแบบหลายแบบพร้อมกันช่วยให้ผู้ใช้ให้ความคิดเห็นได้เปิดกว้างและลึกกว่าการทดสอบการใช้งานจากดีไซน์แบบเดียว~\cite{tohidi2006getting} อีกทั้งการสร้างต้นแบบหลายแบบยังนำไปสู่ผลงานที่ดีขึ้นและมีความหลากหลายของแนวคิดมากขึ้น~\cite{dow2010parallel}
  \end{itemize}
  จากการศึกษาข้างต้น โครงงานนี้เลือกใช้ \textit{Comparative Usability Test}
  เป็นหลัก โดยให้ผู้เข้าร่วมทดสอบเปรียบเทียบ two-alternative UI และให้ข้อเสนอแนะในแต่ละรูปแบบ เนื่องจากโครงงานนี้มุ่งเน้นการค้นหาปัญหาและความต้องการของผู้ใช้จากต้นแบบ เพื่อนำข้อเสนอแนะที่ได้กลับไปปรับปรุงการออกแบบในรอบถัดไป
  ทั้งนี้ Generative AI จะถูกนำมาใช้เพื่อช่วยสร้าง two-alternative UI รวบรวม
  และวิเคราะห์ข้อเสนอแนะ เพื่อลดภาระงานของนักออกแบบในการทดสอบแต่ละรอบ
}

\subsection{การสกัดความต้องการและการสร้าง User Persona จากข้อเสนอแนะของผู้ใช้}
\thaijustify{
  การทำความเข้าใจความต้องการของผู้ใช้เป็นส่วนสำคัญของกระบวนการออกแบบ
  โดยการใช้ต้นแบบเป็นหนึ่งในวิธีที่ช่วยให้นักออกแบบสามารถสกัดความต้องการ
  และรับข้อเสนอแนะจากผู้ใช้ได้ก่อนการพัฒนาระบบจริง \cite{pacheco2018}
  UI prototype  ยังใช้ทดสอบ ประเมิน และยืนยันแนวคิดร่วมกับผู้ใช้และทีมพัฒนาได้~\cite{kolthoff2024}โครงงานนี้เลือกใช้ต้นแบบระดับ high-fidelity เพื่อให้ผู้เข้าร่วมทดสอบสามารถชี้ข้อเสนอแนะไปยัง component ที่ต้องการได้โดยตรง ข้อมูลที่ได้จากการทดสอบต้นแบบจึงสามารถนำมาใช้เพื่อค้นหาปัญหาการใช้งาน ความต้องการ และความแตกต่างระหว่างผู้ใช้แต่ละกลุ่มได้
 
}
\thaijustify{
  อย่างไรก็ตาม ข้อเสนอแนะจากผู้ใช้ไม่ควรถูกตีความว่าเป็นปัญหาการใช้งาน
  หรือความต้องการโดยตรงทั้งหมด เนื่องจากความชอบของผู้ใช้อาจไม่สอดคล้อง
  กับประสิทธิภาพการใช้งานเสมอไป \cite{nngroup_satisfaction} และการใช้ผลการโหวตจากผู้เข้าร่วมจำนวนไม่มากเป็นข้อมูลเชิงปริมาณ อาจทำให้เกิดความเข้าใจผิดได้ \cite{nngroup_qualitative} นอกจากนี้ ข้อเสนอแนะจากผู้มีส่วนได้ส่วนเสียในแต่ละกลุ่มอาจมีความสำคัญแตกต่างกันตามบทบาท ความต้องการ ความเร่งด่วน และอิทธิพลต่อโครงการ \cite{shukla2025,mitchell1997} ดังนั้น การวิเคราะห์ข้อเสนอแนะจึงควรพิจารณาทั้งเนื้อหาของข้อเสนอแนะและบริบทของกลุ่มผู้ให้ข้อเสนอแนะร่วมกัน
}
\thaijustify{
  วิธีหนึ่งในการนำข้อมูลดังกล่าวมาใช้สนับสนุนกระบวนการออกแบบคือ
  User Persona ซึ่งเป็นตัวแทนของกลุ่มผู้ใช้ที่สร้างขึ้นจากข้อมูลของผู้ใช้
  เพื่อสรุปลักษณะและความต้องการที่สำคัญของแต่ละกลุ่ม งานของ
  Shin et al. \cite{shin2024} ศึกษาการใช้ Large Language Models
  ร่วมกับมนุษย์ในการสร้าง Persona จากข้อมูลผู้ใช้ ซึ่งพบว่าการให้ LLM
  สร้าง Persona จากข้อมูลทั้งหมดโดยอัตโนมัติยังมีข้อจำกัดในการระบุลักษณะสำคัญของกลุ่มผู้ใช้ ในทางกลับกัน การให้มนุษย์มีส่วนร่วมในการ
  ระบุและจัดกลุ่มข้อมูลตามลักษณะที่สำคัญ แล้วใช้ LLM ช่วยสรุปข้อมูลของแต่ละกลุ่มเป็น Persona สามารถสร้าง Persona ที่เป็นตัวแทนของข้อมูลผู้ใช้ได้ดีกว่า แนวทางดังกล่าวจึงสนับสนุนแนวคิด Human-in-the-loop
  แทนการให้ Generative AI เป็นผู้กำหนด Persona เพียงอย่างเดียว
}
\thaijustify{
  จากแนวคิดดังกล่าว โครงงานนี้จึงนำข้อเสนอแนะที่ได้จาก Comparative Usability Test มาใช้เป็นข้อมูลเชิงคุณภาพ โดยแยกพิจารณาปัญหาการใช้งาน ความชอบ และกลุ่มของผู้ให้ข้อเสนอแนะ จากนั้นใช้ Generative AI ช่วยรวบรวมและสรุปลักษณะร่วมของข้อมูลเพื่อสนับสนุนการสร้าง User Persona โดยให้นักออกแบบเป็นผู้ตรวจสอบและยืนยันการจัดกลุ่มและผลลัพธ์ที่ได้ Persona ที่ผ่านการตรวจสอบแล้วจึงสามารถใช้เป็นตัวแทนของความต้องการของผู้ใช้แต่ละกลุ่ม และเป็นข้อมูลประกอบการตัดสินใจในการปรับปรุง two-alternative UI ในรอบการออกแบบถัดไปได้
}

\subsection{ข้อกำหนดระบบการออกแบบด้วย Generative AI}
\thaijustify{
  การทำงานกับ Generative AI เพื่อสร้าง two-alternative UI จำเป็นต้องมีกรอบควบคุม (Guardrails) เพื่อให้ผลลัพธ์ที่ได้สอดคล้องกับอัตลักษณ์ของผู้ว่าจ้างหรือผลิตภัณฑ์นั้น ๆ Figma \cite{figma_agent} ระบุว่าการใช้ Design System ช่วยเป็นกรอบให้ Generative AI สามารถสร้างทางเลือกการออกแบบที่ถูกต้องได้ และ Design Tokens Community Group \cite{designtokens2025} ได้กำหนดมาตรฐานรูปแบบไฟล์สำหรับการแลกเปลี่ยน Design tokens ระหว่างเครื่องมือไว้ โครงงานนี้จึงกำหนดให้ Design System เป็นข้อบังคับใช้ควบคุมในการสร้าง UI ด้วย Generative AI และใช้ตรวจสอบความถูกต้องก่อนการส่งมอบ
}
